\documentclass[10pt]{article}
\usepackage[margin=20mm]{geometry}
\usepackage{amsmath,amssymb,mathtools,fontspec,unicode-math,xurl,hyperref}
\IfFontExistsTF{DejaVu Serif}{\setmainfont{DejaVu Serif}}{\setmainfont{Latin Modern Roman}}
\setmathfont{latinmodern-math.otf}
\setlength{\parindent}{0pt}\setlength{\parskip}{6pt}\sloppy
\title{Technical Review of Invertibility of Functionals of the Poisson Process and Applications}
\author{AI-assisted technical review}
\date{}
\begin{document}
\maketitle
Abstract: The acquired version contains an incompatible composition calculus and an invalid construction in the entropy approximation argument. Explicit examples in the bounded, strictly positive class demonstrate these failures. Four additional local issues concern a missing inverse-clock symbol, the admissible compensator domain, the common-noise premise of pathwise uniqueness, and a filtration enlargement. The evidence distinguishes false intermediate assertions from final conclusions: no counterexample to the principal entropy criterion, variational identity, or corrected fixed-driver uniqueness theorem is established.\par
\section*{Identity, scope, and notation}
This report concerns Laure Coutin and Laurent Decreusefond, Invertibility of functionals of the Poisson process and applications, DOI 10.1214/24-AOP1748. The exact source is the 41-page arXiv:2307.02887v2 PDF, submitted 28 November 2024 [paper]. All page locators refer to this PDF and agree with its printed page numbers. Equivalence with the 2025 journal-final text is not asserted.\par
The complete source text, definitions, proofs, and references were read. The review follows the time-change calculus, the invertibility criterion, the entropy approximation, and the weak/strong generalized Hawkes formulation. The paper’s external change-of-measure and martingale representation results were treated as imported dependencies, not independently re-proved from their cited monographs.\par
Let \(N\) be the canonical unit Poisson process under \(\pi\), with natural filtration \(\mathcal N\). Write \[y(t)=\int_0^t\dot y(s)\,ds,\qquad y^*=y^{-1},\qquad
Y(t)=N(y^*(t)),\qquad m(x)=(1+x)\log(1+x)-x.\] The strong class \(P_\infty^{++}\) requires a predictable intensity bounded above and away from zero and equal to one after a fixed finite time [paper, p. 8]. The deterministic test intensities below satisfy this condition wherever the relevant theorem assumes it. Their integrated clocks are finite, strictly increasing homeomorphisms; their inverse times are deterministic stopping times. Their finite-horizon likelihoods are strictly positive, uniformly integrable, and have expectation one, and all displayed entropy costs are finite. Thus the decisive examples do not depend on explosion, degeneracy, or a change-of-measure integrability failure.\par
\section*{Composition calculus}
Equation reference: eq:double. Problem 1 (Major: incompatible process composition).  Definition 2.8 on page 11 defines \[
\mathbf Y(U)(N,t)=U(Y,y^*(N,t)).\] Under this definition, the stochastic-integral identity (24) in Theorem 2.7, page 12, is false. Its proof on page 13 uses an inconsistent stopped-process composition. The same calculus is invoked in the inverse-likelihood argument of Theorem 3.7, page 21, and the discussion and proof of Theorem 3.9, pages 23–27.\par
Decisive example. Set \(\dot y(s)=2\) for \(0<s\leq1\), and \(\dot y(s)=1\) thereafter. Then \[y(t)=\begin{cases}2t&t\leq1,\\t+1&t>1,\end{cases}
\qquad
y^*(t)=\begin{cases}t/2&t\leq2,\\t-1&t>2.\end{cases}\] Choose the deterministic predictable integrand \(r=1\), whose integral is well defined at every finite time. With \(U(N,t)=N(t)\), the left side of (24), interpreted by [eq:double] (\#eq:double), is \(N(t/4)\) for \(0<t\leq2\). The right side is \[\int_0^t r(Y,y^*(s))\,dY(s)=Y(t)=N(t/2).\] Their difference is Poisson with parameter \(t/4\). At \(t=1\), they differ with probability \(1-e^{-1/4}>0\).\par
The failure also survives a compactly supported integrand and a terminal-time reading. For \(r(s)=\mathbf1_{(0,1]}(s)\) and \(t=2\), the left side is \(Y(1)=N(1/2)\), while the right side is \(Y(2)=N(1)\). The discrepancy has mean \(1/2\) and is nonzero with probability \(1-e^{-1/2}\). ◻\par
Alternative readings and counterevidence. Using ordinary configuration pullback \(U(Y,t)\) instead of [eq:double] (\#eq:double) does not validate both printed rules: in the Lebesgue-integral identity (23), \(r=1\) gives \(t\) on the left and \(y^*(t)\) on the right. Pure time substitution \(U(N,y^*(t))\) has its own valid change-of-variable rule, with the integrand evaluated at \((N,y^*(s))\); it is distinct from configuration pullback. Theorem 2.6 supplies the usual time-change rule. These three operations need separate notation and proofs.\par
The example does not disprove the main entropy criterion. Here \(Y\) is invertible, with inverse \(Z(N,t)=N(y(t))\), and the correct entropy is \[H(Y_\#\pi\mid\pi)=\int_0^\infty m(\dot y^*(s)-1)\,ds
=1-\log2.\] For this deterministic clock, the inverse-likelihood identity can be checked directly by ordinary configuration pullback. Consequently the established defect is an invalid core calculus and a resulting proof dependency requiring repair, not a false final entropy theorem. Repair must distinguish configuration pullback, time substitution, and their combination, then recheck the likelihood and entropy deductions for random predictable clocks.\par
\section*{Entropy approximation}
Problem 2 (Major: the approximants in Theorem 4.3 are not valid).  The construction on pages 30–31 does not establish the approximation claimed in Theorem 4.3. There are three independently checkable failures in the displayed intensity and step approximation. Theorem 4.2 explicitly depends on this approximation to obtain its variational representation [paper, p. 29].\par
The theorem assumes a density \(L=d\nu/d\pi\) with \(L\log L\in L^1(\pi)\) and \(\log L\in L^r(\pi)\) for some \(r>1\). It seeks \(\dot y_n\in P_{\infty,\mathrm{pc}}^{++}\) and convergence involving the density \(d(Y_n)_\#\pi/d\pi\).\par
Equation reference: eq:factor. Three decisive tests. First, page 31 prints \[
\dot y(N,s)=\frac{\mathbb E_\pi[D_sL\mid\mathcal N_s]}{\mathbb E_\pi[L\mid\mathcal N_s]},
\qquad D_sF(N)=F(N+\epsilon_s)-F(N).\] The admissible bounded density \(L=1\) has \(D_sL=0\), so [eq:factor] (\#eq:factor) gives zero, whereas its likelihood intensity must be one. If this zero factor is used up to time \(T>0\), the likelihood is \(e^T\mathbf1_{\{N(T)=0\}}\), not the identity density. The usual likelihood factor includes the additive one; with predictable versions it is \[1+\frac{\mathbb E[D_sL\mid\mathcal N_{s-}]}{\mathbb E[L\mid\mathcal N_{s-}]}.\] For \(0<m\leq L\leq M\), its inserted-density form gives bounds \(m/M\) and \(M/m\). This explains both the missing constant and the needed strict positivity.\par
Second, the displayed step construction sets \(\dot y_n=0\) on \([0,T/n)\), then claims \(\dot y_n\in P_{\infty,\mathrm{pc}}^{++}\). Each approximant violates the lower bound on an interval of positive length. Its integrated clock is flat there and is not the required homeomorphism. This is not a value-at-one-point convention.\par
Equation reference: eq:wronggrid. Third, on \([iT/n,(i+1)T/n)\), the printed average is \[
n\int_{(i-1)/n}^{i/n}\dot y(N,s)\,ds,\qquad i=1,\ldots,n-1.\] The integration interval and factor do not match the time mesh unless \(T=1\). To respect also the bounded-density reduction on page 30, take \(T=2\), \(a=1/2\), and set \[g(N,s)=1+\mathbf1_{(0,a]}(s)\mathbf1_{\{N(s-)=0\}},
\qquad \tau=T_1\wedge a.\] This intensity is predictable, lies between one and two, and equals one after \(a\). Its integrated clock is \(y(t)=t+\min(t,\tau)\). Its terminal likelihood satisfies \[L=\Lambda_g=(1+\mathbf1_{\{T_1\leq a\}})e^{-\tau},
\qquad e^{-a}\leq L\leq2,\] and \[\mathbb E_\pi L=e^{-2a}+\int_0^a2e^{-2u}\,du=1.\] Consequently \(L\log L\) and \(\log L\) have all the required integrability, including the proof’s lower and upper density bounds. The clock is a strictly increasing unbounded homeomorphism, its inverse times are stopping times, and \(g\in P_\infty^{++}\).\par
Equation reference: eq:l2. On \(A=\{N(a)=0\}\), which has probability \(e^{-a}\), the intensity is exactly two on \((0,1/2]\) and one afterwards. Let \(g_n\) be the printed construction, including its zero first interval. For \(n\geq4\) divisible by four, direct pathwise integration on \(A\) gives \[
\int_0^2|g_n(s)-g(s)|^2\,ds=\frac12+\frac{10}{n}.\] The zero first interval contributes \(8/n\), while the erroneously prolonged value two on \((1/2,1+2/n)\) contributes \(1/2+2/n\). Therefore \[\mathbb E_\pi\int_0^2|g_n(s)-g(s)|^2\,ds
\geq e^{-1/2}\left(\frac12+\frac{10}{n}\right)
\not\longrightarrow0.\] Even replacing the first interval by the correct positive value leaves the limiting grid error \(1/2\) on \(A\). The retained exact rational grid calculations apply to this pathwise template; the stopped intensity ensures that the counterexample also satisfies the bounded-density reduction. ◻\par
Repair and impact boundaries. A positive initial value and the average \((n/T)\int_{(i-1)T/n}^{iT/n}\dot y(s)\,ds\) correct the immediate mesh defects. They do not by themselves complete the theorem. The proof works with likelihoods \(\Lambda_{y_n}\), whereas the conclusion is stated for densities of \((Y_n)_\#\pi\). Page 23 explicitly warns that these are different objects. For the deterministic clock of Problem Problem 1 (\#p:composition), \(\pi_y\) has rate two on \((0,1]\), while \(Y_\#\pi\) has rate \(1/2\) on \((0,2]\). A justified inverse-clock passage and the required entropy convergence must therefore be supplied.\par
This is a demonstrated failure of the printed construction. It is not a counterexample to the existence statement of Theorem 4.3 or to the final variational identity: for \(L=1\), the identity clock already gives the desired approximation. The severity concerns an essential proof input, not an asserted impossibility of repair.\par
\section*{Additional local problems}
Problem 3 (Minor: the wrong clock in Lemma 3.8).  Equation (41), page 22, bounds \(\mathbb E[-\log\Lambda^*_{y^*}]\) by \(\mathbb E\int_0^\infty m(\dot y(s)-1)\,ds\). Under the stated hypotheses this bound is false; the inverse-clock derivative is needed.\par
Take \(\dot y=a=1/2\) on \((0,1]\), and one afterwards. This belongs to \(P_2^{++}(\mathcal N,\pi_y)\); \(\pi_y\) and \(\pi\) are equivalent, \(\mathbb E_\pi\Lambda_y=1\), and the assumed inverse-clock cost is finite. Since \(\Lambda_y=a^{N(1)}e^{1-a}=\Lambda^*_{y^*}\) at terminal time, \[\begin{aligned}
\mathbb E_\pi[-\log\Lambda^*_{y^*}]&=\log2-\tfrac12,\\
\int_0^\infty m(\dot y-1)\,ds&=\tfrac12(1-\log2).\end{aligned}\] The left side exceeds the printed right side by \(\tfrac32\log2-1>0\). Strict positivity follows, for example, from \(\log2=2\sum_{k\geq0}3^{-(2k+1)}/(2k+1)>2/3\). The correct inverse intensity is two until time \(1/2\), so its cost is \(\tfrac12(2\log2-1)=\log2-\tfrac12\).\par
The proof on page 23 and the main entropy criterion already use the inverse-clock cost; equation (30), page 16, gives the correct change of variables. These are substantial counterevidence against a defect in the final criterion arising from this symbol alone. Replacing \(\dot y\) by \(\dot y^*\) repairs the displayed bound in the tested setting. The severity is local and minor, separately from Problem Problem 1 (\#p:composition).\par
Problem 4 (Minor: overbroad compensator existence).  Theorem 2.1(2), page 6, asserts existence of a probability law on the simple configuration space for every predictable, nondecreasing, right-continuous process \(y\), null at zero, with \(N-y\) a local martingale localized by the jump times. That domain is too broad.\par
Set \(y(t)=t+2\mathbf1_{[1,\infty)}(t)\). It is deterministic and predictable, right-continuous, nondecreasing, finite at finite times, null at zero, and unbounded. On the simple configuration space of Definition 2.1, \(\Delta N(1)\in\{0,1\}\). A compensator must satisfy, at this predictable time, \[\Delta y(1)=\mathbb E[\Delta N(1)\mid\mathcal N_{1-}]\leq1,\] in contradiction with \(\Delta y(1)=2\). Localization gives the same necessary jump identity, so allowing a local martingale does not remove the contradiction.\par
An admissible compensator domain and appropriate existence/nonexplosion conditions must be specified. A jump bound alone should not be presented as a full general existence theorem. The later principal arguments mostly use absolutely continuous clocks and explicitly assume an existing law with stronger regularity. This example therefore refutes the preliminary universal statement, not those later claims under their full hypotheses; that limits its severity to a local domain error.\par
Problem 5 (Minor: pathwise uniqueness lacks a common-noise premise).  Definition 5.5, page 33, requires any two solutions \((Z,N)\) and \((Z',N')\) on one filtered probability space to satisfy both \(Z=Z'\) and \(N=N'\), without assuming the equality of the drivers.\par
For \(\dot y=1\), put two independent unit Poisson processes \(N_1,N_2\) on their product space with the joint right-continuous filtration. The pairs \((N_1,N_1)\) and \((N_2,N_2)\) both satisfy every item of Definition 5.2: they are locally finite simple point processes, \(y(Z,t)=t\) is a stopping time, each compensated driver is a local martingale, and the defining equation holds. Nevertheless \[\mathbb P\{N_1(1)=0,\ N_2(1)=1\}=e^{-2}>0,\] so the conclusion of the definition fails even for the identity clock.\par
The repair is to assume \(N=N'\) as the common driver and conclude indistinguishability of \(Z,Z'\). The proof of Theorem 5.5 on page 37 already compares solutions driven by one \(N\), and the coupling in Theorem 5.3 uses a common third coordinate. Those passages support the intended interpretation. The result is a defective definition and literal formulation, not a disproof of fixed-driver uniqueness or of a corrected Yamada–Watanabe statement.\par
Problem 6 (Minor: unjustified enlargement of a martingale filtration).  The proof of Lemma 5.2, page 35, asserts that a square-integrable martingale in the driver’s natural filtration is consequently a martingale in the same-time product filtration \(\mathcal N^{\otimes(2)}\). This implication is false for the allowed solution pairs.\par
Use the bounded deterministic clock of Problem Problem 1 (\#p:composition) and the valid solution \(Z(t)=N(y(t))\), on the original natural filtration of \(N\). Let \(s=1/4\), \(t=1/2\), and \(F=\mathbf1_{\{N(t)=0\}}\). The bounded natural-filtration martingale \[M_r=\mathbb E[F\mid\mathcal N_r]-\mathbb E[F]\] has the square-integrable Poisson representation used in the proof. In the product filtration \(\mathcal G_s=\sigma(Z(u),N(u):u\leq s)\), the value \(Z(s)=N(1/2)=N(t)\) is already known. Hence \(M_t=F-e^{-t}\) is \(\mathcal G_s\)-measurable. On \(\{N(t)=0\}\), an event of probability \(e^{-1/2}\), \[\mathbb E[M_t\mid\mathcal G_s]-M_s=1-e^{-1/4}>0.\] Thus \(M\) is not a martingale in the stated enlarged filtration. All solution assumptions hold because the clock times are deterministic stopping times in the original filtration; Definition 5.2 does not require \(Z(t)\) to be measurable at physical time \(t\) there.\par
The preceding part of the proof claims that the relevant event \(B\) is measurable in the driver’s \(\mathcal N_t\). Given that assertion, the cancellation in the second step can use the original martingale property directly. The time-changed filtrations in the subsequent coupling also differ from this same-time product filtration. The false enlargement assertion is therefore removable from that cancellation, and no additional failure of the final theorem is inferred from it. Repair should delete it or state and establish a suitable compatibility condition.\par
\section*{Overall assessment and reviewed-version limits}
Problems Problem 1 (\#p:composition) and Problem 2 (\#p:approx) establish invalid steps in central arguments. Problems Problem 3 (\#p:cost)–Problem 6 (\#p:filtration) are more limited formulation or local proof errors, with explicit counterevidence bounding their impact. The review does not equate any of these with whole-paper failure. The last section of the paper expressly distinguishes its approach from the preceding entropy construction; no blanket transfer of the entropy proof defects to every Hawkes conclusion is warranted.\par
The primary executor checked the official arXiv metadata on 29 October 2026; it identifies v2 dated 28 November 2024. A bounded official-source search did not identify a separate formal erratum. The author’s institutional publications page was unavailable to the primary open request; a helper’s metadata result is retained separately. Neither search establishes absence of corrections or equality with the journal-final version. No author code was executed. Missing optional code is not treated as a mathematical defect. This is a completed initial technical review; scientific adjudication remains pending and human reviews are empty.\par
\section*{Supporting checks}
`5a6de8b8d509c3cc3bc97de9c065c76e9d2dd8edd02f503017ba0e21563afa26`.\par
The per-problem records retain locators, hypotheses, analytic proofs, alternative interpretations, dependencies, and repair boundaries. The reviewer-written `independent\_checks.py` records exact rational and symbolic calculations, including the pathwise template in [eq:l2] (\#eq:l2); it executes no author material. The bounded-density refinement and its positive expected-error lower bound are retained in `evidence/P2\_bounded\_density\_check.json`, with the likelihood normalization checked symbolically. Source evidence images, actual inspection records, the exact TeX and compiled PDF, compiler dependency payloads, logs, and current inventories are retained in this paper’s own verification tree. Native helpers performed file integrity and metadata work only. Product UI observations and runtime self-report are separately attributed, without a claim of backend equivalence.\par
9 L. Coutin and L. Decreusefond, Invertibility of functionals of the Poisson process and applications, arXiv:2307.02887v2 (28 November 2024). https://arxiv.org/abs/2307.02887v2. Journal DOI: https://doi.org/10.1214/24-AOP1748.\par
\end{document}
